% 40_body.tex — 산출물 ㊵ 편익 실현 원장 공동 확정본 (빈 템플릿 / 완성 예시 공용 본문) — gen_product_templates.py 가 Product_specs.py 에서 생성
% 드라이버: 40_편익원장공동확정본_빈템플릿.tex (\wbblanktrue) / 40_편익원장공동확정본_완성예시.tex (\wbblankfalse — 워크북 §X.5 확정 후)
\documentclass[10pt,a4paper]{article}
\usepackage[a4paper,margin=16mm,top=14mm,bottom=20mm]{geometry}
\def\DocNum{40}\def\DocDay{8}\def\DocIdx{5}
\def\DocTitle{편익 실현 원장 공동 확정본}
\def\DocSub{Benefits Realisation Ledger — Joint Confirmation (= PM ⑳)}
\def\DocVariant{\B{빈 템플릿}{딥게이지 완성 예시}}
\usepackage{wbstyle}
\def\DocCirc{㊵}   % newunicodechar(㉛~㊵ 폴백) 뒤에 정의해야 활성 문자로 읽힌다
\begin{document}
\docheader
 {\Bg{[온담 P1 후속 · 딥게이지 GAUGE]}{온담 P1 후속 '검사기록 디지털화' · 딥게이지 GAUGE}}
 {\Bg{[v2.0 / 2028-05]}{v2.0 / 2028-05-31 (D+821)}}
 {\Bg{[작성자 — 양측 PM]}{나영선(온담) · 윤하경(딥게이지)\rd{7mm}}}
 {\Bg{[서명자 — 나영선 · 강민수 · 측정 책임자]}{나영선 · 강민수 · 한동석 · 오정근 · 조성태\rd{7mm}}}

%% ------------------------------------------------------------------ 1
\secbar{1. 문서 정보와 당사자}
\guide{PM ⑳ 편익 원장 v1.0(2027-05-21, 재무본부 소유)을 상위 문서로. 딥게이지가 output 열을 채우고 양측이 서명하는 공동 확정본.}
{\footnotesize
\begin{tabularx}{\linewidth}{|L{36mm}|Y|}\hline
\hd{항목} & \hd{내용} \\ \hline
\ifwbblank
\rd{8mm}상위 문서(⑳ 버전·소유) &  \\ \hline
\rd{8mm}당사자·서명자 &  \\ \hline
\rd{8mm}기준일 / 갱신 시점 &  \\ \hline
\rd{8mm}적용 범위(편익 항목) &  \\ \hline
\else
상위 문서(⑳ 버전·소유) & 편익 실현 원장 v1.0(2027-05-21, G4 종료 보고서 부속, 소유 재무본부 정미란·재경팀장). 이 문서 = ⑳ 리뷰 계획의 '2028-05 운영 +12 원장 v2 — B-02 분모 B 목표 설정'. 소유는 재무본부, 딥게이지는 output 열 서명자 \\ \hline
당사자·서명자 & output — 딥게이지 강민수 / outcome — 한동석(B-02) · 오정근(B-03·04) · 나영선(B-09·10) / utilisation — 나영선 · 점포 QA · 강현도 / 조성태(안성 세트) / 판정 — 재무본부 \\ \hline
기준일 / 갱신 시점 & 2028-05-31 · G5(2028-03-31) 결과 반영 · 다음 갱신 2029-05(운영 +24) 및 온담 세트 측정 확정 시(서명 + 8주) \\ \hline
적용 범위(편익 항목) & B-02(분모 B) · B-03 · B-04 + 신설 B-09 · B-10. B-02(A)·05·07·08은 ⑳ 그대로(참고) \\ \hline
\fi
\end{tabularx}}

%% ------------------------------------------------------------------ 2
\landscapeon
\secbar{2. 편익 항목 — ⑳의 6칸 + output}
\guide{온담 B-01~08 중 검사기록 디지털화가 기여하는 항목(B-02 재고 정확도 분모 B, B-03 결품, B-04 폐기). 측정식·기준값(분모 명시)·목표값·측정 시점·측정 책임자 실명·미달 시 조치 + **딥게이지 output**(검사기록 자동 생성률·코드 매핑 커버리지·판정 정확도).}
{\footnotesize
\begin{tabularx}{\linewidth}{|L{18mm}|Y|L{34mm}|L{34mm}|L{30mm}|Y|}\hline
\hd{편익 ID} & \hd{측정식(분자/분모)} & \hd{기준값·시점} & \hd{목표값·시점} & \hd{측정 책임자(실명)} & \hd{딥게이지 output} \\ \hline
\ifwbblank
\rd{11mm} &  &  &  &  &  \\ \hline
\rd{11mm} &  &  &  &  &  \\ \hline
\rd{11mm} &  &  &  &  &  \\ \hline
\rd{11mm} &  &  &  &  &  \\ \hline
\rd{11mm} &  &  &  &  &  \\ \hline
\rd{11mm} &  &  &  &  &  \\ \hline
\else
B-02(A) 참고 & 실사 일치 SKU ÷ 순환실사 SKU(69\%) & A 91.3\%(2025) → G5 97.1\% & 98.5\%(M+24 2028-01, 판정 완료 — 81\% 실현) & 한동석 / 순환실사 담당 & (검사기록과 무관 — ⑳ 그대로) \\ \hline
B-02(B) 분모 B & 실사 일치 SKU ÷ 전체 SKU(100\%, 비순환 31\% 포함) & B 84.6\%(2026-12-18) → G5 88.9\% & 95.0\%(M+36 2030-01) — v2 신설 & 한동석 / 나영선(검사기록 전제) & 입고 검수 검사기록 자동 생성률 ≥ 90\% · 코드 매핑 커버리지 61.4 → 95\% · 판정 정확도: 온담 세트 미측정(항목 4) \\ \hline
B-03 결품률 & 결품 SKU-일 ÷ 진열 SKU-일 & 4.7\%(2025) → G5 2.6\% & 1.5\%(M+24) — 66\% 실현, 유지 판정 2029-05 & 오정근 / 매장운영 KPI 담당 & 입고 검수 리드타임(검사 대기 재고) 1.8일 → 0.5일 \\ \hline
B-04 폐기율 & 월 폐기 금액 ÷ 취급원가(C-15) · 수량 병기 & 3.9\%(2025) → G5 2.9\% & 2.2\%(M+24) — 59\% 실현, 개선 계획 & 오정근 / 강현도 & '온도 이탈' 폐기 코드 자동 부여율 ≥ 95\%(C-27) · out-of-range → 폐기 전환율 기준 31\% \\ \hline
B-09 검사기록 작성·검색 공수(신설) & 월 작성 인시 + 검색 인시(61점 + 이천 + 안성) & 월 1,260 인시(2028-04 표본) & −60\% = 504 인시(운영 +12 2029-11) & 나영선 / 점포 QA 담당 & 기록 자동 생성률 · 검색 응답 3초 · 오프라인 동기화 성공률 \\ \hline
B-10 CCP 초과 → 개선조치 기록 리드타임(신설) & 한계 초과 시각 → 개선조치 기록 완료 시각(61점 평균) & 3.2시간(2028-04 표본) & 30분(운영 +12) & 나영선 / 오정근(점포 실행) & 플래그 → 알림 30초(C-03) · 판정은 검사원 — 법정 기준 판정 정확도는 output 아님(㊲ 항목 5) \\ \hline
\fi
\end{tabularx}}
\landscapeoff

%% ------------------------------------------------------------------ 3
\secbar{3. 책무 분리 — output / outcome / utilisation}
\guide{output(딥게이지 — 무엇을 인도하는가) / outcome(온담 현업 임원 — 무엇이 달라지는가) / utilisation(업무 변화 소유자 — 누가 쓰게 만드는가). 칸마다 실명.}
{\footnotesize
\begin{tabularx}{\linewidth}{|L{18mm}|Y|Y|Y|}\hline
\hd{편익 ID} & \hd{output — 딥게이지} & \hd{outcome — 온담 임원} & \hd{utilisation — 업무 변화 소유자} \\ \hline
\ifwbblank
\rd{11mm} &  &  &  \\ \hline
\rd{11mm} &  &  &  \\ \hline
\rd{11mm} &  &  &  \\ \hline
\rd{11mm} &  &  &  \\ \hline
\rd{11mm} &  &  &  \\ \hline
\else
B-02(B) & 딥게이지 강민수 — 자동 생성률·커버리지·측정 세트 80건 인도 & 한동석 물류본부장 — 분모 B 95\% & 나영선 — 검사기록 정착 · 마스터 정비(C-05, 온담 자체) \\ \hline
B-03 & 딥게이지 — 검사 대기 재고 리드타임 리포트 & 오정근 영업본부장 & 매장운영 KPI 담당 · 구매(입고 검수) \\ \hline
B-04 & 딥게이지 — 온도 이탈 코드 자동 부여(C-27) & 오정근 & 강현도 구매팀장 · 한동석(냉장 운영) \\ \hline
B-09 & 딥게이지 — 모바일·오프라인·검색 & 나영선 & 점포 QA · 지역매니저(체크리스트 C-44) \\ \hline
B-10 & 딥게이지 — CCP 플래그·알림(C-03) & 나영선 & 점포 검사원(개선조치 실행) · 오정근 \\ \hline
\fi
\end{tabularx}}

%% ------------------------------------------------------------------ 4
\secbar{4. 인수 기준의 모집단}
\guide{FN 2.4\% · OCR 95.2\%는 자동차 3세트 60건의 값이다. 온담 4세트(냉장·HACCP·입고·안성)에서의 측정 계획(세트 구성·건수·시점·판정자)을 인수 기준 조항으로. 모집단 없는 fit criterion은 무효.}
{\footnotesize
\begin{tabularx}{\linewidth}{|L{40mm}|Y|}\hline
\hd{항목} & \hd{내용} \\ \hline
\ifwbblank
\rd{10mm}현재 기준의 모집단(V-07) &  \\ \hline
\rd{10mm}온담 4세트 측정 계획 &  \\ \hline
\rd{10mm}임시 인수 기준(측정 전)과 유효 기간 &  \\ \hline
\rd{10mm}측정 후 확정 절차·판정자 &  \\ \hline
\rd{10mm}측정 비용의 부담 &  \\ \hline
\else
현재 기준의 모집단(V-07) & 자동차 3세트(A·B·C) 60건 — 오전 조도, 절삭·프레스·사출. FN 2.4 / OCR 95.2 / FP 21.3은 이 60건의 값. 온담 4세트에서는 한 번도 측정되지 않았다. 기준은 두 번 개정(㉗ FN ≤ 8.0 → ㉞ 3.0), 데이터셋은 불변 \\ \hline
온담 4세트 측정 계획 & D-1 이천 냉장 20(한동석 → 나영선) / D-2 61점 HACCP 20(나영선) / D-3 입고 성적서 상위 30개사 20(나영선 + 구매) / D-4 안성 납품 20(조성태 → 나영선) = 80건. 서명 후 8주(라벨링 4 + 측정 2 + 검토 2). 라벨 판정자 = 정본 판정자(㊲ 항목 7) \\ \hline
임시 인수 기준(측정 전)과 유효 기간 & D-3·D-4(건 단위): OCR ≥ 90\% \& FN ≤ 5\% / D-2: '개선조치 필요' 재현율 ≥ 97\% / D-1: in-range 판정 정확도 ≥ 99\%(OCR 없음). 측정 전 자동판정 OFF·초안만. 유효: 측정 확정 시까지, 최장 서명 + 12주 \\ \hline
측정 후 확정 절차·판정자 & 측정 후 4주 내 양측 서명 — 나영선(온담) · 오세진(딥게이지). 미확정 시 임시 기준 유지 + 초안 OFF 지속. 확정 기준 = ㊲ 항목 4 'AI 판정 정확도' 행의 목표 \\ \hline
측정 비용의 부담 & 1,200만 — 딥게이지 부담(㊳ 사용처 개발·eval 행). 인수 기준을 만드는 비용은 공급자의 것 — 온담에 청구하면 측정이 미뤄진다 \\ \hline
\fi
\end{tabularx}}

%% ------------------------------------------------------------------ 5
\secbar{5. 측정 시점·리뷰 — 두 시계}
\guide{온담 G5(2028-03-31은 이미 지났다 — 후속 리뷰 시점) · M+12/M+24 · 딥게이지 갱신 시점(분기 릴리스). 어느 리뷰에 누가 오는가.}
{\footnotesize
\begin{tabularx}{\linewidth}{|L{34mm}|L{26mm}|Y|Y|}\hline
\hd{리뷰} & \hd{시점} & \hd{참석·판정} & \hd{판정 대상} \\ \hline
\ifwbblank
\rd{10mm} &  &  &  \\ \hline
\rd{10mm} &  &  &  \\ \hline
\rd{10mm} &  &  &  \\ \hline
\rd{10mm} &  &  &  \\ \hline
\else
온담 G5 후속 — 원장 v2 확정 & 2028-05-31(운영 +12) & 재무본부(정미란·재경팀장) 판정 · 나영선 · 강민수 배석 & B-02 분모 B 목표 신설 · B-09·10 신설 · 측정 계획 조항 \\ \hline
온담 4세트 측정 확정 & 계약 서명 + 8주 & 나영선(판정자) · 오세진 · 조성태 & 임시 인수 기준 → 확정 기준 \\ \hline
운영 +24 유지 판정 & 2029-05 & 재무본부 · 운영위 · 딥게이지 배석(판정자 아님) & B-03·04·08 유지 · B-09·10 첫 판정 · 원장 종결 여부 \\ \hline
딥게이지 분기 릴리스 리뷰 & 분기(온담 승인 후) & 박세연 · 딥게이지 온담 담당 · 나영선 위임자 & output 지표 3종 · 오류 예산 · 미감시 라인 \\ \hline
\fi
\end{tabularx}}

%% ------------------------------------------------------------------ 6
\secbar{6. 미달 시 조치·배상의 분리}
\guide{편익 미달(outcome)과 output 미달(계약 불이행)을 분리한다. output 미달만 ㊲ 배상 조항의 대상.}
{\footnotesize
\begin{tabularx}{\linewidth}{|L{40mm}|L{30mm}|Y|L{26mm}|}\hline
\hd{상황} & \hd{책임 주체} & \hd{조치} & \hd{배상 연동 여부} \\ \hline
\ifwbblank
\rd{10mm}output 미달 &  &  &  \\ \hline
\rd{10mm}outcome 미달(output 충족) &  &  &  \\ \hline
\rd{10mm}utilisation 미달 &  &  &  \\ \hline
\rd{10mm}측정 불가 &  &  &  \\ \hline
\else
output 미달 & 딥게이지 & 재이관·재측정 · 무상 기간 연장 & 연동 — ㊲ 25\% 구간 \\ \hline
outcome 미달(output 충족) & 온담 현업 임원(한동석·오정근·나영선) & 개선 계획(⑳ 규칙) — 딥게이지는 데이터 제공 & 비연동 \\ \hline
utilisation 미달 & 업무 변화 소유자(나영선·점포 QA) & 교육·체크리스트·지역매니저 점검 & 비연동 \\ \hline
측정 불가 & 온담 물류본부(분모 B 실사 자체가 없음) & 측정 체계 먼저 — 판정 보류 & 비연동 — 단 측정 세트 80건은 output \\ \hline
\fi
\end{tabularx}}

%% ------------------------------------------------------------------ 7
\secbar{7. 8일 사슬 역추적}
\guide{이 표의 각 숫자·이름이 어느 문서의 어느 행에서 왔는가 — ⑳ ← ⑱ ← ① / ㊵ ← ㉗ ← ㉖ ← ㉓ ← ㉒.}
{\footnotesize
\begin{tabularx}{\linewidth}{|Y|Y|Y|}\hline
\hd{이 문서의 칸} & \hd{출처 문서·항목} & \hd{변경 이력} \\ \hline
\ifwbblank
\rd{9mm} &  &  \\ \hline
\rd{9mm} &  &  \\ \hline
\rd{9mm} &  &  \\ \hline
\rd{9mm} &  &  \\ \hline
\rd{9mm} &  &  \\ \hline
\rd{9mm} &  &  \\ \hline
\else
항목 2 B-02 분모 A 91.3 / B 84.6 & ⑳ 항목 2 분모 정치(2027-05-21) ← ⑫ Day3 분모 정정(2026-09) ← ① BC §4 '순환실사 SKU'(2025-11) & 분모 A는 3년째 그대로 — B 목표만 신설 \\ \hline
항목 3 output / outcome / utilisation 실명 & ⑳ 항목 5 ITO 서명표 ← ④ 거버넌스·RACI(Day1) & 딥게이지가 output 열 서명자로 추가 \\ \hline
항목 4 '자동차 3세트 60건' & ㉞ 항목 5(D+540) ← ㉗ 항목 1(D+137) ← ㉖ 항목 3 범위 '자동차'(D+130) ← ㉓ 초기 SAM 2,150(D+60) ← ㉒ 박선재 관찰(D+40) & 세트 D는 한 번도 없었다 — 이 칸이 비면 서명 불가 \\ \hline
항목 4 FN 2.4\% · 임시 기준 & ㉞ eval v2 합격 기준 개정(D+540, criteria drift) ← ㉗ 항목 3 개정 절차 · v1 FN ≤ 8.0 & 기준은 두 번 바뀌었고 모집단은 한 번도 안 바뀌었다 \\ \hline
항목 6 output 미달만 배상 & ㊲ 항목 5 4구간 ← ㉞ 항목 4 162만(D+525) ← ㉘ 항목 6 책임 한계 10\%(D+137) ← ⑰ SLA 배상 조항(Day4, 요구하는 쪽) & 같은 조항을 반대편에서 \\ \hline
항목 5 두 시계 & ⑳ 항목 4 리뷰 계획 ← ⑱ 이관 6번(2027-05) ← 2025-11 이사회 반려 34억 SI 견적 · ② 헌장 §5 제외 범위 & '제외 범위'가 3년 뒤 21억 계약이 됐다 \\ \hline
\fi
\end{tabularx}}

%% ------------------------------------------------------------------ 8
\secbar{8. 양측 서명}
\guide{output — 딥게이지 / outcome·utilisation — 온담. 서명 전 확인: 인수 기준의 모집단 칸이 채워졌는가.}
{\footnotesize
\begin{tabularx}{\linewidth}{|Y|Y|Y|Y|}\hline
\hd{딥게이지 강민수 CEO} & \hd{온담 나영선 품질안전본부장} & \hd{측정 책임자(실명)} & \hd{온담푸드시스템 조성태} \\ \hline
\ifwbblank
\rd{16mm} & & &  \\ \hline
\else
\rd{16mm} 강민수 CEO (서명) & 나영선 품질안전본부장 (서명) & 한동석 · 오정근 · 나영선 (서명) & 조성태 대표 (서명) — 항목 4 채움 확인 \\ \hline
\fi
\rd{7mm}일자 & & &  \\ \hline
\end{tabularx}}

\end{document}